% §7 Evaluation — target 2.25 pages
% Structure mirrors RQ1-RQ4.

\section{Evaluation}
\label{sec:eval}

% ── §7.1 overall ──────────────────────────────────────────────────────────
\subsection{Overall performance (RQ1)}
\label{sec:eval:overall}

\begin{table*}[t]
\centering
\caption{Main results — Recall / Precision / F1 / weighted Score, averaged over
7 scenarios $\times$ 3 runs per model. Weighted Score normalizes by
$\sum w(g)$ for $g \in \text{GT}$.}
\label{tab:main}
\footnotesize
\begin{tabular}{@{}lrrrrr@{}}
\toprule
Model & Recall & Precision & F1 & Score & Cost (\$) \\
\midrule
Claude Sonnet 4     & \todo{} & \todo{} & \todo{} & \todo{} & \todo{} \\
Gemini 2.5          & \todo{} & \todo{} & \todo{} & \todo{} & \todo{} \\
GPT-4.1             & \todo{} & \todo{} & \todo{} & \todo{} & \todo{} \\
DeepSeek-V3         & \todo{} & \todo{} & \todo{} & \todo{} & \todo{} \\
MiniMax/Qwen3       & \todo{} & \todo{} & \todo{} & \todo{} & \todo{} \\
\bottomrule
\end{tabular}
\end{table*}

\todo{Write §7.1 prose — headline finding: best LLM reaches XX\% recall at /24
scale, variance across runs stays under YY\%.}

% ── §7.2 baselines ────────────────────────────────────────────────────────
\subsection{Network-scale vs baselines (RQ4)}
\label{sec:eval:baselines}

\begin{figure}[t]
  \centering
  % \includegraphics[width=\linewidth]{figures/fig4_baselines.pdf}
  \fbox{\parbox{0.95\linewidth}{\centering\textit{Fig.~4 placeholder: bar chart,
  F1 and Score for \harnessname vs OpenVAS, Nmap NSE, CAI-adapted,
  VulnBot-adapted.}}}
  \caption{Baseline comparison. Scanners capture known CVEs but miss
  configuration vulnerabilities. Single-host LLM agents lose cross-device
  signal when adapted to a network scope.}
  \label{fig:baselines}
\end{figure}

\todo{Write §7.2 — explicit narrative: (i) scanners win on CVE category,
(ii) LLMs win on misconfig/default-creds, (iii) single-host LLMs adapted to
network fall behind native harness on F1 by ZZ\%.}

% ── §7.3 per-architecture ─────────────────────────────────────────────────
\subsection{Per-architecture analysis (RQ3)}
\label{sec:eval:arch}

\begin{figure}[t]
  \centering
  % \includegraphics[width=\linewidth]{figures/fig5_heatmap.pdf}
  \fbox{\parbox{0.95\linewidth}{\centering\textit{Fig.~5 placeholder:
  heatmap scenario $\times$ model, cell = weighted Score.}}}
  \caption{Per-scenario Score heatmap. \todo{Note which pattern consistently
  breaks LLMs.}}
  \label{fig:heatmap}
\end{figure}

\todo{Write §7.3 — pattern-by-pattern commentary. Expected finding:
gateway-centric and hub-star degrade most, because they require pivot
reasoning. Edge-cloud may also degrade on cloud-side service discovery.}

% ── §7.4 evidence level depth ─────────────────────────────────────────────
\subsection{Evidence-level depth (RQ2)}
\label{sec:eval:evidence}

\begin{table}[t]
\centering
\caption{Evidence-level coverage. $L_k$ is the fraction of true positives
reaching level~$k$ or deeper. Absent from binary-flag benchmarks.}
\label{tab:evidence}
\footnotesize
\begin{tabular}{@{}lrrr@{}}
\toprule
Model & $L_1$ (detect) & $L_2$ (exploit) & $L_3$ (exfiltrate) \\
\midrule
Claude Sonnet 4  & \todo{} & \todo{} & \todo{} \\
Gemini 2.5       & \todo{} & \todo{} & \todo{} \\
GPT-4.1          & \todo{} & \todo{} & \todo{} \\
DeepSeek-V3      & \todo{} & \todo{} & \todo{} \\
MiniMax/Qwen3    & \todo{} & \todo{} & \todo{} \\
\midrule
OpenVAS          & --      & --      & --      \\
Nmap NSE         & --      & --      & --      \\
\bottomrule
\end{tabular}
\end{table}

\todo{Write §7.4 — expected: detection (L1) near-parity with scanners, but
L2/L3 is LLM-only. This is the strongest argument for LLM agents in this
context.}

% ── §7.5 per-protocol ─────────────────────────────────────────────────────
\subsection{Per-protocol analysis}
\label{sec:eval:protocol}

\todo{Write §7.5 — 2 paragraphs. Expected: LLMs excel on MQTT/SSH/HTTP
(text-rich, well-represented in training), fall short on Modbus/LoRaWAN
(binary, sparse training data). Report numbers per protocol.}

% ── §7.6 failure modes, cost, time ────────────────────────────────────────
\subsection{Failure modes and cost/time envelope}
\label{sec:eval:failures}

\begin{figure}[t]
  \centering
  % \includegraphics[width=\linewidth]{figures/fig6_pareto.pdf}
  \fbox{\parbox{0.95\linewidth}{\centering\textit{Fig.~6 placeholder:
  Pareto plot, USD cost vs weighted Score, one point per model.}}}
  \caption{Cost--performance Pareto. \todo{Annotate which model dominates
  on each axis and the efficient frontier.}}
  \label{fig:pareto}
\end{figure}

\todo{Write §7.6 — categorize failure modes: (a) hallucinated findings on
devices with no services, (b) missed vulnerabilities on binary protocols,
(c) over-reporting info-disclosure without severity calibration. Brief
quantitative table.}
